\documentclass[11pt]{letter}
\usepackage[margin=1in]{geometry}
\usepackage[hidelinks]{hyperref}
\signature{Pushkal Gupta\\Corresponding author, on behalf of all authors\\
School of Computer Science and Engineering\\
Vellore Institute of Technology, Vellore 632014, Tamil Nadu, India\\
\href{mailto:pushkal.gupta2023@vitstudent.ac.in}{pushkal.gupta2023@vitstudent.ac.in}}
\address{}
\date{\today}
\begin{document}
\begin{letter}{The Editor\\Applied Energy}
\opening{Dear Editor,}

We submit the manuscript \emph{``What does heterogeneity cost? A controlled, communication-priced decomposition of
federated LSTM load forecasting on London smart meters''} by Pushkal Gupta, Vanka Kaushik and Kalaavathi B for
consideration as a Research Article in \emph{Applied Energy}.

Federated learning is increasingly proposed as the privacy-preserving way to forecast household electricity demand
from smart-meter data. Two of its premises are rarely measured: that its accuracy loss is caused by heterogeneity
between households, and that it saves bandwidth compared with collecting the data. Using 100 households from the
public Low Carbon London trial, we hold the homes, the test set, the model and the compute budget fixed and change
one factor at a time, while counting every byte exchanged. The main findings are:
\begin{itemize}
\item The federated protocol on a single client matches centralised accuracy. The accuracy gap is caused by
fragmentation: it grows with the number of clients the same homes are split over, from 2\,\% with five clients to
10\,\% with one client per home.
\item Heterogeneity between clients, dialled over real load-profile clusters and measured with a label-free index,
barely changes final accuracy, but it raises the traffic needed to approach centralised accuracy by about 85\,\%.
\item With a small LSTM and half-hourly data, federated training moves more bytes than uploading the raw readings.
The case for federation is therefore governance, not bandwidth.
\item Plain FedAvg leaves most households worse off than local models. One epoch of local fine-tuning, at no
communication cost, largely removes this at community scale. Proximal regularisation does not help, and we show why
an update-dissimilarity-adaptive variant fails.
\end{itemize}
These results give distribution system operators, suppliers and aggregators concrete guidance: federate a few
aggregated clients rather than individual meters, fine-tune locally, and budget heterogeneity as communication. We
believe this fits the journal's focus on data-driven methods for energy systems with clear practical implications.
All code, configurations and per-run results are provided, and every number in the manuscript is generated from them.

This manuscript is original, has not been published, and is not under consideration elsewhere. All authors have
approved the manuscript and agree with its submission. The authors declare no competing interests. The use of
AI-assisted tools in preparing the software and the text is disclosed in the manuscript.

Thank you for considering our work.

\closing{Yours sincerely,}
\end{letter}
\end{document}
